\section{Symmetry, decoration data, and the scope of the theory}\label{sec:scope}
The purpose of this reference is to make a complete decorated state, its obstruction, and its stacking law readable as one mathematical object. The central problem is the integer Chern-insulator layer. Charge conservation, fermionic exchange signs, and the phase generated by removing an exact decoration must be retained together. A nonzero obstruction cochain is not an obstruction to a state unless its class survives all permitted lower-layer changes.

The formulas below use one convention: the calibrated CI/charge representative selected by the source comparison and the bordism product test in Sections~\ref{sec:normalization} and~\ref{sec:examples}. There are no version-labelled alternatives in the operative formulas. The cochain operations, dimensional equations, equivalences, and calibration arguments have separate roles and separate locations.

\paragraph{Scope.}
The reference constructs and calibrates the CI-generated, localized-charge, and group-cohomological phase sector through $4+1$ dimensions. It retains the proofs for the specified natural lower-coordinate model. This is not a claim that the independent neutral $E_8$ sector, every intrinsic higher-dimensional response, or every universal parity comparison map has been constructed. The exact boundaries are collected in Section~\ref{sec:boundaries}; they are not additional free coefficients in the adopted $4+1$D formula.

\subsection{The same fermionic symmetry and its two quotients}
The spin--charge relation and the two quotients are
\begin{equation}\label{eq:quotients}
 e^{i\pi Q}=(-1)^F,\qquad
 G=G_f/\Uf,\qquad G_b=G_f/\mathbb Z_2^f,
 \qquad BU(1)_b\longrightarrow BG_b\xrightarrow{\pi}BG.
\end{equation}
Here $U(1)_b=\Uf/\mathbb Z_2^f$. Insulator cochains are on $G$; the parity cochains used for comparison are on $G_b$. These are two filtrations of the same symmetry problem, not two unrelated symmetry groups. Keeping the quotient names distinct is necessary even though both are sometimes called the ``bosonic group.''

The symmetry characters are defined on operators:
\begin{equation}\label{eq:characters}
 gQg^{-1}=(-1)^{\rho_1(g)}Q,\qquad
 gig^{-1}=(-1)^{s_1(g)}i,\qquad
 g e^{i\theta Q}g^{-1}=e^{i(-1)^{\rho_1(g)+s_1(g)}\theta Q}.
\end{equation}
Thus ordinary charge-preserving time reversal has $\rho_1=0$, $s_1\ne0$. The circle automorphism is $\rho_1+s_1$, not $\rho_1$. In Ref.~\cite{WQFCG} the character denoted by $\rho_1$ is the circle automorphism; it is the sum of the two characters used here.

The extension and its explicit real lift are
\begin{equation}\label{eq:omega}
 \omega_2\in Z^2\bigl(G,\mathbb R/\mathbb Z_{\rho_1+s_1}\bigr),
 \qquad \omega_2^{\mathbb R}\bmod1=\omega_2,
 \qquad d_{\rho_1+s_1}\omega_2^{\mathbb R}\in C^3\bigl(G,\mathbb Z_{\rho_1+s_1}\bigr).
\end{equation}
The binary extension in the parity comparison is denoted by $\omega_2^{\SC}$, not by an unqualified $w_2$. The notation $w_i(TM)$ is reserved for Stiefel--Whitney classes in the geometric examples.

\subsection{Fields, outputs, and one notation for each object}
A subscript always gives the cochain degree. Each physical dimension has its own declaration:
\begin{center}
\begin{tabular}{llll}\toprule
Spacetime & CI number & Localized charge & Additive phase\\\midrule
$2+1$ & $n_0\in C^0(G,\mathbb Z_{s_1})$ & $n_2\in C^2(G,\mathbb Z_{\rho_1})$ & $\nub_3\in C^3(G,\mathbb R/\mathbb Z_{s_1})$\\
$3+1$ & $n_1\in C^1(G,\mathbb Z_{s_1})$ & $n_3\in C^3(G,\mathbb Z_{\rho_1})$ & $\nub_4\in C^4(G,\mathbb R/\mathbb Z_{s_1})$\\
$4+1$ & $n_2\in C^2(G,\mathbb Z_{s_1})$ & $n_4\in C^4(G,\mathbb Z_{\rho_1})$ & $\nub_5\in C^5(G,\mathbb R/\mathbb Z_{s_1})$\\\bottomrule
\end{tabular}
\end{center}
The equations in the dimensional sections impose closure and the higher constraints; this table does not declare every field closed.

Primes mark a second input, and capital letters mark the output: $n'_j$, $N_j$, $\mathcal V_j$. Following the canonical AHSS formula guide~\cite{AHSSGuide}, obstruction functions and stacking twisters are $\mathcal O_j$ and $\mathcal E_j$. Hats distinguish additive phases:
\begin{equation}\label{eq:phase-convention}
 \nu_j=e^{2\pi i\nub_j},\quad
 \mathcal V_j=e^{2\pi i\Vb_j},\quad
 \mathcal O_j=e^{2\pi i\Ob_j},\quad
 \mathcal E_j=e^{2\pi i\Eb_j}.
\end{equation}
The last two exponentials apply only to phase-valued operations. An unhatted integer $\mathcal O_j$ or $\mathcal E_j$ is an integer cochain, not a phase. On individual modified fields the prime remains outside the modifier: $\bar n'_j$, $\widetilde n'_j$, $\check n'_j$.

\subsection{The one derived background needed by the CI formulas}
The antiunitary correction is packaged by
\begin{equation}\label{eq:AB}
 \boxed{\displaystyle
 \check\omega_2^{\mathbb R}
 =\omega_2^{\mathbb R}-\frac12(s_1\cup s_1)_{\rho_1+s_1},
 \qquad B_3=d_{\rho_1+s_1}\check\omega_2^{\mathbb R}.}
\end{equation}
Both quantities have coefficient sign $\rho_1+s_1$, and $B_3$ is integral. The array $s_1\cup s_1$ is first formed as an ordinary integer cocycle from the canonical bits; the subscript then assigns the stated coefficient type. This shift does not change the physical symmetry extension. It is a declared coordinate construction whose effects are included in the decoration equations and phases.

The superscript $\mathbb R$ is essential: $\check\omega_2^{\mathbb R}$ is a real cochain, whereas $\check\omega_2^{\SC}=\omega_2^{\SC}+s_1^2$ in the parity comparison is binary. No bare-letter substitutes for $s_1$, $\rho_1$, or their sum are used in reader formulas.

\subsection{Which neutral data are outside this sector?}
The full two-dimensional electronic invertible coefficient has two response coordinates~\cite{BCHM}:
\begin{equation}\label{eq:neutral-coefficient}
 (n_0,e_0)\in\mathbb Z^2,\qquad
 \sigma_{xy}=n_0,\qquad c_-=n_0+8e_0,\qquad
 n_0^{\SC}=2n_0+16e_0.
\end{equation}
Here $\sigma_{xy}$ is measured in $e^2/h$, and $e_0$ counts neutral $E_8$ layers. Parity-only theory already contains $E_8$ as sixteen $p+ip$ layers; adding it again there would double count. Retaining charge conservation distinguishes electric from thermal response, so the forgetful map is not injective. For example $(8,-1)$ has zero thermal response and nonzero electric response.

The present equations set the independent neutral decoration aside; they do not assign its obstruction or product the value zero. Likewise, a relative $2+1$D SPT calculation sets both uniform integers to zero, but this does not prohibit a neutral sheet decoration in higher dimension. Throughout the examples, a CI-generated subgroup is distinguished from the full bordism classification.
